\section{Introduction}\label{sec:intro}

\subsection{Exactly divergence-free elements}
Most mixed finite element methods for incompressible flow enforce the constraint $\div u=0$ only weakly: the discrete divergence is orthogonal to the discrete pressure space, but not zero. In 1985 Scott and Vogelius \cite{ScottVogelius1985} observed that for continuous piecewise-$P_k$ velocities and discontinuous piecewise-$P_{k-1}$ pressures the divergence maps the velocity space \emph{onto} the pressure space when $k\ge4$ and the mesh has no singular vertices. For such a pair the weak constraint is a strong one: the discrete velocity is divergence-free pointwise. The inf-sup constant of the pair was later made explicit in terms of the local geometry of the mesh by Guzm\'an and Scott \cite{GuzmanScott2019}, through an angle quantity $\Theta(z)$ at each vertex $z$ that measures how far the vertex is from being singular. Lower-degree pairs with the same property exist on split meshes, for instance on barycentric (Clough--Tocher) refinements \cite{ArnoldQin1992,GuzmanNeilan2018}; see Neilan's review \cite{Neilan2020}.

Pointwise incompressibility is more than a cosmetic property. It gives exact mass conservation, and, on meshes fitted to the domain with strongly imposed Dirichlet data, it makes the velocity error independent of the pressure. This \emph{pressure robustness} matters when pressure gradients are large relative to the viscosity, for example in nearly hydrostatic or rotating flows; the survey of John, Linke, Merdon, Neilan and Rebholz \cite{JohnEtAl2017} explains why the classical inf-sup stable pairs lack it and why divergence-free pairs have it.

\subsection{Curved walls and polygonal approximation}
Most flows of interest are bounded by curved walls. The simplest way to handle a curved boundary $\Gamma$ in a finite element code is to replace it by an inscribed polygon $\Gh$ whose vertices lie on $\Gamma$, and to impose the boundary condition on $\Gh$. For scalar second-order problems the consequences are classical. Berger, Scott and Strang \cite{BergerScottStrang1972} and Scott \cite{Scott1975} showed that transferring a homogeneous Dirichlet condition from $\Gamma$ to $\Gh$ costs an $H^1$ error of order $\hG^{3/2}$, where $\hG$ is the length of the chords: the solution is of size $\hG^2$ on the chords and oscillates on the scale $\hG$, so its $H^{1/2}(\Gh)$ norm is of order $\hG^{3/2}$. Higher order requires either a better geometry or a correction of the boundary data. The first route is the isoparametric approach, in which the boundary elements are curved by a polynomial map \cite{Ciarlet1978book}; the second goes back to Bramble, Dupont and Thom\'ee \cite{BrambleDupontThomee1972}, who corrected the boundary condition on the polygon by a Taylor expansion in the normal direction. Both ideas are alive today: boundary-value correction has been combined with Nitsche's method for cut elements \cite{BurmanHansboLarson2018,BurmanHansboLarson2019Stokes} and for polygonal approximations \cite{DupontGuzmanScott2025}, and the shifted boundary method \cite{MainScovazzi2018} is a related approach.

For exactly divergence-free elements the polygonal approximation has an additional, less obvious cost. A field that is smooth at a polygon vertex and vanishes on both adjacent edges has zero gradient there, since its gradient annihilates two independent directions. Continuous piecewise polynomials can escape this constraint, because their gradients may jump across the interior edges at the vertex; for \emph{divergence-free} piecewise polynomials, however, the jumps are themselves constrained, and such fields inherit gradient constraints at the vertices of a polygonal wall \cite{ScottPrize,GjerdeScott2024}. On the true curved wall, however, the gradient of the velocity is $\dn u\otimes n$, which is not zero where the wall shear stress is not zero. If no-slip is imposed strongly on an inscribed polygon, the computed gradient can therefore be wrong by $O(1)$ at the polygon vertices. Scott and Gjerde--Scott report that the $H^1$ error of Scott--Vogelius elements then decays only like $h^{1/2}$ \cite{ScottPrize,GjerdeScott2024}, compared with $h^{3/2}$ for scalar problems. We call this the \emph{zero-gradient lock}. The standard cure is to fit the geometry, with isoparametric or Piola-mapped Scott--Vogelius elements \cite{NeilanOtus2021,DurstNeilan2024}; these attain the optimal order $h^k$, at the price of curved elements.

\subsection{Nitsche's method and the prize question}
In 1971 Nitsche \cite{Nitsche1971} proposed to impose Dirichlet conditions weakly, by adding to the variational form a consistency term (the normal flux of the solution tested against the test function), its symmetric counterpart, and a penalty term $(\mu/h)\ip{u}{v}$. Unlike a pure penalty method \cite{Babuska1973}, whose penalty must grow as the mesh is refined because the method is inconsistent, Nitsche's method is consistent, so the penalty only has to be large enough for coercivity. Stenberg \cite{Stenberg1995} later showed that Nitsche's method can be read as a stabilised Lagrange multiplier method, which explains both its consistency and the form of its penalty; the multiplier is the boundary flux of the solution. For the Stokes equations the consistency term is the full traction $\sigma(u,p)n=(\nabla u)n-pn$ (or its symmetric-gradient version), and therefore contains the pressure term $\ip{p}{v\cdot n}$ \cite{FreundStenberg1995,Stenberg1995,Becker2002,BurmanHansbo2014}; Burman, Hansbo and Larson note that this term is essential for consistency \cite{BurmanHansboLarson2019Stokes}.

Gjerde and Scott \cite{GjerdeScott2024} combined these ideas: they keep the inscribed polygon, use Scott--Vogelius elements with $k=4$, and impose no-slip on the polygon by Nitsche's method, with penalty $\mu/h$ and \emph{no pressure term}. They report errors orders of magnitude below those of strong imposition. Scott observed $H^1$ errors of order $h^{3/2}$ for this Scott--Vogelius--Nitsche method on a benchmark with constant pressure, and posed the question as a prize problem, listed as PPL~115 in the Prize Problem Ledger \cite{ScottPrize}. The prize asks for ``a published proof that Scott--Vogelius--Nitsche converges'', which ``should include a computational study of the influence of the Nitsche penalty parameter''. It states the expected error
\[
\norm{u-u_h}_{H^1}\;\lesssim\;\hG^{3/2}+\hO^k,
\]
where $\hG$ is the boundary mesh size and $\hO$ the bulk mesh size, and it adds: \emph{``If this does not hold, why not?''} This paper answers the question, for general Stokes data, in two dimensions.

\subsection{What goes wrong, and why the method still converges}
The key observation is simple. Because the Gjerde--Scott form has no pressure term on the wall, the discrete momentum equation does not see the wall pressure through a traction. Testing the exact solution in the discrete equations leaves the residual $-\ip{p-\pbar}{v\cdot n}$ on the polygon, for any constant $\pbar$. Something must balance it, and the only candidate is the penalty. To leading order the discrete velocity therefore satisfies the \emph{leak law}
\[
\frac\mu h\,u_h\cdot n\;\approx\;p-\pbar\qquad\text{on }\Gh ,
\]
where $\pbar$ is the mean of $p$ over the wall: the fluid leaks through the wall, outward where the pressure is above its mean and inward where it is below. The leak is of size $h/\mu$, so it vanishes as the mesh is refined. This is the classical consistency error of a boundary penalty method \cite{Babuska1973}, acting only on the pressure and with the vanishing parameter $h/\mu$, and it is the reason why the printed method converges at all.

How fast does it converge? The leak is a normal velocity of size $h/\mu$ on the wall, so a naive trace argument would lift it to an $H^1$ error of order $(h/\mu)\hG^{-1/2}$, that is, rate $\frac12$. The true $H^1$ rate is $1$. The reason is incompressibility: for a divergence-free polynomial field on a boundary triangle, the normal derivative of the normal component equals minus the tangential derivative of the tangential component, so the purely \emph{normal} load produced by the missing traction acts like a tangential derivative and can be integrated by parts along the wall. In the mesh-dependent energy norm, which contains $(\mu/h)^{1/2}\norm{\cdot}_{L^2(\Gh)}$, the leak is seen directly and the rate is exactly $\frac12$. In short: the printed method converges, but not at Scott's rate, unless the pressure is constant along the wall. This explains why Scott's constant-pressure benchmark shows the rate $\frac32$, and it is the answer to ``why not'' for the method as printed.

\subsection{The repair}
The obvious repair is to add the missing term $\ip{p_h}{v\cdot n}$ to the momentum equation. With exactly divergence-free elements this does not work as it stands: the square linear system acquires the kernel vector $(u,p)=(0,1)$, because a constant pressure now drops out of both equations. The standard fix, restricting the pressure to mean-zero functions, destroys pointwise incompressibility: the discrete divergence becomes a nonzero constant. The repair we analyse adds the traction with the \emph{wall mean of the discrete pressure removed}, $\ip{p_h-\pbG(p_h)}{v\cdot n}$, to the momentum equation only. The continuity equation is untouched, so the velocity stays exactly divergence-free, and the constant mode is fixed automatically. We call this method $\CNS$ (consistent Nitsche, with the star marking the mean-free correction). The device is not new: Frachon, Nilsson and Zahedi \cite{FrachonNilssonZahedi2024} introduced the zero-net-flux formulation for divergence-free cut elements, and the velocity of $\CNS$ coincides with the velocity of their formulation transcribed to our setting. What is new is the analysis on an inscribed polygon.

\subsection{Main results}
We state the results informally here; precise statements are in Sections~\ref{sec:printed}--\ref{sec:strong}. We write $h=\hO$ for the bulk mesh size (the $h$ in Scott's penalty $\mu/h$), $\hG$ for the longest chord, $\gamma=\mu\hG/h$ for the effective boundary penalty, and $G=\norm{p-\pbar}_{L^2(\Gamma)}$ for the variation of the pressure along the wall. The solution is assumed smooth, the mesh shape regular, and every polygon vertex is assumed to lie in at least three triangles (this is the Guzm\'an--Scott angle condition at the wall; Section~\ref{sec:setting}).
\begin{enumerate}[label=(\arabic*)]
\item \emph{The printed method converges, but at rate $1$ when the wall pressure varies} (Section~\ref{sec:printed}). Its $H^1$ error is at most $C_p\,h/\mu+C(\hG^{3/2}+\hO^k)$, and at least $c\,(h/\mu)G$; in the energy norm it is $\asymp(h/\mu)^{1/2}G$ plus Scott's terms. To leading order the error is an explicit field whose trace is $-(h/\mu)(p-\pbar)n$; its slip and energy norms, normalised by $(h/\mu)G$ and $(h/\mu)^{1/2}G$, tend to $1$. The normalised $H^1$ error tends to the energy of the Stokes flow that has the leak as Dirichlet data. This constant depends on the whole domain, not only on the wall.
\item \emph{The mean-free correction attains Scott's rate, and the rate is sharp} (Section~\ref{sec:corrected}). $\CNS$ is well posed for $\gamma$ above a threshold, its velocity is divergence-free pointwise, and $\norm{u-u_h}_{H^1}\le C(\hG^{3/2}+\hO^k)$ for general Stokes data and every fixed $k\ge4$. With a zero-flux correction of the outer boundary data the bound is uniform in the penalty and needs no relation between $\hG$ and $\hO$. If the wall shear does not vanish identically and $\hG\le2\sin(\theta_{\min}/4)$, with $\theta_{\min}$ the minimum angle of the triangles at the wall, the $H^1$ error is also bounded below by $c\,\hG^{3/2}$, so the term $\hG^{3/2}$ cannot be improved.
\item \emph{The penalty must scale with the mesh ratio} (Section~\ref{sec:penalty}). With Scott's normalisation $\mu/h$, the penalty must grow like $\rho=h/\min_e|e|$. Sufficiency holds on every mesh. Necessity is certified in exact rational arithmetic on every mesh whose shortest chord carries a triangle with apex angle at least $35^\circ$ and base angles at least $5^\circ$, and is proved on every shape-regular mesh with explicit but very small constants. A necessity statement based on a single wall star is false in general (the exact counterexamples have small angles).
\item \emph{The zero-gradient lock is a property of the wall mesh} (Section~\ref{sec:strong}). Under the same angle condition, strongly imposed no-slip on the polygon already has $H^1$ error $\asymp\hG^{3/2}+h^k$, two-sided. The lock is caused by polygon vertices lying in only two triangles, which the angle condition excludes: each such vertex costs $\hG|\dn u(z)|$ in $H^1$, one lock gives rate $1$, and a positive fraction of locks gives rate $\frac12$. With three or more triangles the same loss appears continuously as one triangle at the vertex degenerates. The advantages of Nitsche imposition on an inscribed polygon are therefore robustness to such meshes and a smaller constant, not a better rate.
\end{enumerate}
Table~\ref{tab:summary} summarises the rates.

\begin{table}[htb]
\centering\small
\caption{Rates for $k\ge4$ on meshes satisfying the angle condition (M0)--(M2) of Section~\ref{sec:setting}, fixed penalty $\gamma$, smooth solution, $h\asymp\hG$. ``$G>0$'' means the pressure is not constant on the wall; ``$W\not\equiv0$'' means the wall shear does not vanish identically.}\label{tab:summary}
\setlength{\tabcolsep}{3pt}
\begin{tabular}{>{\raggedright\arraybackslash}p{0.29\textwidth}>{\raggedright\arraybackslash}p{0.33\textwidth}>{\raggedright\arraybackslash}p{0.15\textwidth}>{\raggedright\arraybackslash}p{0.14\textwidth}}
\toprule
method & $H^1$ error & energy error & where\\
\midrule
printed ($\GS$), $G=0$ & $\asymp\hG^{3/2}$ (if $W\not\equiv0$) & $\lesssim\hG^{3/2}$ & Thms.~\ref{thm:A}, \ref{thm:lb}\\
printed ($\GS$), $G>0$ & $\asymp h/\mu$ & $\asymp(h/\mu)^{1/2}$ & Thms.~\ref{thm:A}--\ref{thm:B}\\
printed, growing penalty $\mu\gtrsim h\hG^{-3/2}$ & $\asymp\hG^{3/2}$ & rate $\le1$ for any scaling & Cor.~\ref{cor:rescue}, Sec.~\ref{sec:growing}\\
corrected ($\CNS$), any $G$ & $\asymp\hG^{3/2}$ (if $W\not\equiv0$) & $\lesssim\hG^{3/2}$ & Thms.~\ref{thm:D}, \ref{thm:lb}\\
strong no-slip, $\ge3$ triangles at every wall vertex & $\asymp\hG^{3/2}$ & -- & Thm.~\ref{thm:st:two}\\
strong no-slip, locked wall vertices $\Lk$ & $\asymp\hG^{3/2}+\hG\bigl(\sum_{\Lk}|\dn u(z)|^2\bigr)^{1/2}$ & -- & Thm.~\ref{thm:st:locks}\\
\bottomrule
\end{tabular}
\end{table}

Beyond these four results, the extended version \cite{PadhiExt} of this paper, of which the present article is a condensed account, contains a theory of the pressure and of drag and lift, $L^2$ estimates, Clough--Tocher refinements with $k\ge2$, general curved domains, steady Navier--Stokes flow, nearly singular wall stars, and pressure robustness. We state these results without proof in Section~\ref{sec:further}. The three-dimensional analysis is the subject of a companion paper \cite{PadhiThreeD}.

\subsection{How to read this paper}
Section~\ref{sec:related} compares our results with the recent literature and lists precisely what is new. Section~\ref{sec:setting} fixes the setting and collects the geometric and discrete tools. Section~\ref{sec:printed} analyses the printed method, Section~\ref{sec:corrected} the corrected method, Section~\ref{sec:penalty} the penalty threshold and Section~\ref{sec:strong} strong imposition. Section~\ref{sec:further} surveys further results, Section~\ref{sec:numerics} reports computations, and Section~\ref{sec:conclusions} concludes with open problems. Proofs that are routine, or that rely on long computer-assisted certificates, are sketched and referred to \cite{PadhiExt}; the code and all raw data are publicly available \cite{ZeroGradRepo}.
